\section{Keeping the coefficient and slack budgets together}
\label{mol:sec:joint-budget}

We will replace each atom coefficient by a multiple of the head test's
sum on that segment. This produces two nonnegative errors: the squared
change of coefficients and the unused squared head-test energy. Their
sum is controlled by the squared-radius deficit of the center. Keeping
this joint budget proves the main estimate for the part of an arbitrary
displacement outside the head.

We use the segment atoms, ancestral projections, and cone decomposition
already established above.  The argument applies both to the coordinate
completion on the infinite tree and to the finite-height forest, after
their respective atomic-ball identifications.  It does not identify those
two spaces.



\begin{proof}[Proof of Theorem~\ref{thm:main-tail}]
Put $r=\|x\|$. If $r=0$, contractivity gives
$\|Q_Hy\|\le\|y\|\le R$, which proves the assertion. Assume $r>0$.
First suppose that $y$ has finite support, and enlarge $H$ to a finite
ancestral set $K$ containing that support.  Choose a head test $p$ of
norm one such that $x(p)=r$.  By the finite atomic-ball formula, choose
finite laws of represented vectors
\[
 A=\sum_i a_i d_{S_i},\qquad \sum_i a_i^2\le R^2,
\]
with means $x+y$ and $x-y$.  Choose a sign $\sigma$ with equal
probabilities and then the corresponding law.  Thus
$\mathbb EA=x$ and $\mathbb E(\sigma A)=y$.

For each outcome put
\[
 b_i=r p(S_i),\quad
 A^\circ=\sum_i b_i d_{S_i},\quad
 V=\sum_i(a_i-b_i)^2,\quad
 L=r^2-\sum_i b_i^2.
\]
Both $V$ and $L$ are nonnegative.  Since
$\mathbb E\sum_i a_i b_i=r\mathbb EA(p)=r^2$, expansion gives
\begin{equation}\label{mol:eq:joint-identity}
 \eta+\theta
 :=\mathbb EV+\mathbb EL
 =\mathbb E\sum_i a_i^2-r^2
 \le R^2-r^2.
\end{equation}
The atom coefficient estimate also gives
$\|Q_H(A-A^\circ)\|\le\sqrt V$.

Let $G$ be the gates outside $H$ which occur in the finite
representations.  At a gate $g$, write $b(g)$ for the coefficient of the
crossing segment in $A^\circ$, and put it equal to zero when there is
no crossing segment.  Every segment wholly outside $H$ has replaced
coefficient zero.  If $W=Q_H(A-A^\circ)$, then
\[
 \sum_{g\in G}|W(e_g)|^2\le V.
\]
Indeed, the gates are incomparable, so each segment contains at most
one gate, and disjointness assigns different gates to different
coefficients.  Since $\mathbb EA(e_g)=x(e_g)=0$,
\begin{equation}\label{mol:eq:joint-signed-mass}
 \sum_g(\mathbb Eb(g))^2
 =\sum_g(\mathbb EW(e_g))^2\le\mathbb EV=\eta.
\end{equation}

The simultaneous opposite-head splitting argument, applied to the
test $rp$, gives for two independent outcomes
\[
 2\sum_{g:b(g)b'(g)<0}|b(g)b'(g)|\le L+L'.
\]
Here the available squared-test budget is $r^2$; this is the only
change from the norm-one version.  Consequently, with
$m(g)=\mathbb E|b(g)|$,
\begin{equation}\label{mol:eq:joint-absolute-mass}
 \|m\|_2^2-\|\mathbb Eb\|_2^2
 =\mathbb E\sum_g\bigl(|b(g)b'(g)|-b(g)b'(g)\bigr)
 \le2\theta.
\end{equation}
Equations \eqref{mol:eq:joint-signed-mass} and
\eqref{mol:eq:joint-absolute-mass} therefore give
$\|m\|_2\le\sqrt{\eta+2\theta}$.

On the cone rooted at $g$, each replaced outcome is $b(g)d_{J(g)}$,
where $J(g)$ is its crossing tail starting at $g$.  The term is zero
when there is no crossing segment.  Thus the signed mean decomposes as
\begin{equation}\label{mol:eq:signed-product-recovery}
 Q_H\mathbb E(\sigma A^\circ)
 =\sum_g\mathbb E\bigl[\sigma b(g)d_{J(g)}\bigr].
\end{equation}
Each cone term has norm at most $\mathbb E|b(g)|=m(g)$, since
$|\sigma|=1$ and a segment indicator has norm one.  The cone identity
therefore gives
\[
 \|Q_H\mathbb E(\sigma A^\circ)\|
 \le\|m\|_2\le\sqrt{\eta+2\theta}.
\]
Lemma~\ref{stab:coef:tail-assembly} gives a direct convex-product proof
of the same reconstruction, allowing any multiplier of absolute value
at most one in place of $\sigma$.

The signed mean of the original atoms is $y$, so its tail satisfies
\[
 \|Q_Hy\|\le\sqrt\eta+\sqrt{\eta+2\theta}
 \le2\sqrt{\eta+\theta}.
\]
For the last inequality, square both sides and use
$\sqrt{\eta(\eta+2\theta)}\le\eta+\theta$.
Now apply \eqref{mol:eq:joint-identity}.

For an arbitrary completed vector $y$, take an increasing finite ancestral
exhaustion $K_n\supset H$.  Contractivity gives
$\|x\pm P_{K_n}y\|\le R$, and $P_{K_n}y\to y$ in norm.
The finite estimate passes to the limit.  No bound on the height or
the number of children entered the calculation.
\end{proof}
